% ===========================================================================
% Category: Ensemble Learning
% Generated from ML_Lab_Manual_Completed.md - edit the Markdown and re-run
% scripts/build_latex_manual.py instead of editing this file by hand.
% ===========================================================================

\chapter{Experiment 8: Random Forest Regression
(RFR)}\label{experiment-8-random-forest-regression-rfr}

\textbf{Category:} Ensemble learning

\section{1. Objective}\label{objective}

Predict a continuous target using an ensemble of randomized decision
trees; evaluate with MAE, RMSE and R² and inspect out-of-bag validation.

\section{2. Dataset Used}\label{dataset-used}

California Housing
(\passthrough{\lstinline!fetch\_california\_housing!}), 5,000-row random
sample, 8 features, target = median house value. 80/20 train/test split
(4,000/1,000).

\section{3. Required Libraries}\label{required-libraries}

scikit-learn, numpy, pandas

\section{4. Brief Theory}\label{brief-theory}

Random Forest Regression averages predictions of B trees trained on
bootstrap samples, each split considering a random feature subset.
Averaging decorrelated trees reduces variance. Out-of-bag samples (not
drawn in a bootstrap) give a free internal validation estimate.

\section{5. Procedure}\label{procedure}

\begin{enumerate}
\def\labelenumi{\arabic{enumi}.}
\tightlist
\item
  Load and sample the dataset; split train/test.
\item
  Train
  \passthrough{\lstinline!RandomForestRegressor(n\_estimators=150, max\_depth=12, oob\_score=True)!}.
\item
  Predict the test set; compute R², RMSE, MAE.
\item
  Report the OOB R² and compare with the test R² to check overfitting.
\end{enumerate}

\section{6. Reference Implementation}\label{reference-implementation}

\begin{lstlisting}[language=Python]
rf_reg = RandomForestRegressor(n_estimators=150, max_depth=12,
                               oob_score=True, random_state=42, n_jobs=-1)
rf_reg.fit(X_train_reg, y_train_reg)
print(rf_reg.oob_score_, r2_score(y_test_reg, rf_reg.predict(X_test_reg)))
\end{lstlisting}

\section{7. Results}\label{results}

{\def\LTcaptype{none} % do not increment counter
\begin{longtable}[]{@{}ll@{}}
\toprule\noalign{}
Metric & Value \\
\midrule\noalign{}
\endhead
\bottomrule\noalign{}
\endlastfoot
OOB R² & 0.7513 \\
Test R² & 0.7415 \\
Test RMSE & 0.5908 (\$100k units ≈ \$59k) \\
Test MAE & 0.3999 (\$100k units ≈ \$40k) \\
\end{longtable}
}

OOB R² (0.7513) is very close to the test R² (0.7415), so the model is
not overfitting. Feature importances are dominated by median income and
location (latitude/longitude), consistent with housing economics.

\section{8. Discussion and Conclusion}\label{discussion-and-conclusion}

RFR provides a strong, low-maintenance baseline for tabular regression:
no feature scaling is required, it captures non-linearities and
interactions, and OOB gives free validation. Gradient boosting
(Experiments 10-12) reaches higher R² on the same data at the cost of
more tuning.

\section{9. Answers to Viva Questions}\label{answers-to-viva-questions}

\begin{itemize}
\tightlist
\item
  \textbf{Why does a forest reduce variance?} Averaging B approximately
  uncorrelated trees gives Var ≈ ρσ² + (1−ρ)σ²/B; as B grows the second
  term vanishes, so predictions are more stable than a single tree's.
\item
  \textbf{What is bagging?} Bootstrap aggregating: train each model on a
  bootstrap sample of the data and combine predictions (average for
  regression, vote for classification).
\item
  \textbf{Why can feature importance be misleading?} Impurity-based
  importance is biased toward high-cardinality/continuous features, is
  computed on training data (no causality), and correlated features
  share/dilute importance arbitrarily.
\end{itemize}

\begin{center}\rule{0.5\linewidth}{0.5pt}\end{center}

\section{10. Complete Code Listing}\label{complete-code-listing}

\emph{Full runnable program for this experiment, extracted from
\passthrough{\lstinline!notebooks/04\_ensemble\_learning.ipynb!}
(executed; results above are its actual output).}

\textbf{Listing 8.1 --- Core libraries}

\begin{lstlisting}[language=Python]
# ---- Core libraries ----
import numpy as np
import pandas as pd
import matplotlib.pyplot as plt
import seaborn as sns

from sklearn.datasets import fetch_california_housing
from sklearn.model_selection import train_test_split
from sklearn.ensemble import RandomForestRegressor, RandomForestClassifier, AdaBoostRegressor, AdaBoostClassifier
from sklearn.metrics import r2_score, root_mean_squared_error, mean_absolute_error, accuracy_score, f1_score, roc_auc_score

import xgboost as xgb
import catboost as cb

# Styling setup
plt.style.use('seaborn-v0_8-whitegrid' if 'seaborn-v0_8-whitegrid' in plt.style.available else 'default')
plt.rcParams['font.sans-serif'] = 'DejaVu Sans'
plt.rcParams['figure.figsize'] = (10, 6)
plt.rcParams['figure.dpi'] = 110
np.random.seed(42)
\end{lstlisting}

\textbf{Listing 8.2 --- Load the California Housing dataset (real-world
regression benchmark)}

\begin{lstlisting}[language=Python]
# ---- Load the California Housing dataset (real-world regression benchmark) ----
housing = fetch_california_housing(as_frame=True)
# Sample 5,000 rows so every ensemble trains quickly while remaining representative
df_reg = housing.frame.sample(n=5000, random_state=42)

X_reg = df_reg[housing.feature_names]   # 8 socio-economic / geographic features
y_reg = df_reg['MedHouseVal']           # target: median house value

# ---- Train/test split (80% / 20%) ----
X_train_reg, X_test_reg, y_train_reg, y_test_reg = train_test_split(
    X_reg, y_reg, test_size=0.20, random_state=42
)

print(f"Regression Training Samples: {X_train_reg.shape[0]}, Test Samples: {X_test_reg.shape[0]}")
df_reg.head()
\end{lstlisting}

\textbf{Listing 8.3 --- Random Forest Regression (bagging of decision
trees)}

\begin{lstlisting}[language=Python]
# ---- Random Forest Regression (bagging of decision trees) ----
# n_estimators = number of trees, oob_score = estimate accuracy from out-of-bag samples
rf_reg = RandomForestRegressor(n_estimators=150, max_depth=12, oob_score=True, random_state=42, n_jobs=-1)
rf_reg.fit(X_train_reg, y_train_reg)

# ---- Predict and evaluate ----
y_pred_rf = rf_reg.predict(X_test_reg)
r2_rf = r2_score(y_test_reg, y_pred_rf)                          # explained variance (1.0 = perfect)
rmse_rf = root_mean_squared_error(y_test_reg, y_pred_rf)         # error in target units
mae_rf = mean_absolute_error(y_test_reg, y_pred_rf)              # robust average error

print(f"Random Forest Regressor Results:")
print(f" - Out-of-Bag (OOB) R² Score: {rf_reg.oob_score_:.4f}")
print(f" - Test R² Score:              {r2_rf:.4f}")
print(f" - Test RMSE:                  {rmse_rf:.4f}")
print(f" - Test MAE:                   {mae_rf:.4f}")
\end{lstlisting}

\chapter{Experiment 9: Random Forest Classification
(RFC)}\label{experiment-9-random-forest-classification-rfc}

\textbf{Category:} Ensemble learning

\section{1. Objective}\label{objective-1}

Classify observations with an ensemble of randomized decision trees
using a stratified split and report accuracy, precision, recall, F1 and
the confusion matrix.

\section{2. Dataset Used}\label{dataset-used-1}

Heart Disease (\passthrough{\lstinline!data/heart\_disease.csv!}, 303
rows, 13 features plus target). Categorical variables are one-hot
encoded for the forest; stratified 75/25 split (227/76).

\section{3. Required Libraries}\label{required-libraries-1}

scikit-learn, numpy, pandas

\section{4. Brief Theory}\label{brief-theory-1}

Random Forest Classification aggregates tree votes: Ĉ(x) = argmax\_k
Σ\_b I(T\_b(x) = k). Bootstrap sampling plus feature subsampling
decorrelates trees, reducing variance relative to a single decision
tree. The class-probability output (fraction of votes) enables ROC-AUC.

\section{5. Procedure}\label{procedure-1}

\begin{enumerate}
\def\labelenumi{\arabic{enumi}.}
\tightlist
\item
  Load the heart data; one-hot encode categorical features.
\item
  Use a stratified 75/25 split to preserve the disease ratio.
\item
  Train
  \passthrough{\lstinline!RandomForestClassifier(n\_estimators=120, max\_depth=8)!}.
\item
  Predict labels and probabilities on the test set.
\item
  Compute accuracy, F1 and ROC-AUC.
\end{enumerate}

\section{6. Reference Implementation}\label{reference-implementation-1}

\begin{lstlisting}[language=Python]
rfc = RandomForestClassifier(n_estimators=120, max_depth=8, random_state=42, n_jobs=-1)
rfc.fit(X_train_clf, y_train_clf)
y_prob_rfc = rfc.predict_proba(X_test_clf)[:, 1]
\end{lstlisting}

\section{7. Results}\label{results-1}

{\def\LTcaptype{none} % do not increment counter
\begin{longtable}[]{@{}
  >{\raggedright\arraybackslash}p{(\linewidth - 2\tabcolsep) * \real{0.5000}}
  >{\raggedright\arraybackslash}p{(\linewidth - 2\tabcolsep) * \real{0.5000}}@{}}
\toprule\noalign{}
\begin{minipage}[b]{\linewidth}\raggedright
Metric
\end{minipage} & \begin{minipage}[b]{\linewidth}\raggedright
Value
\end{minipage} \\
\midrule\noalign{}
\endhead
\bottomrule\noalign{}
\endlastfoot
Accuracy & 0.7763 \\
Precision / Recall / F1 & 0.8046 (F1) \\
ROC-AUC & 0.8704 \\
Confusion matrix (test, 76 samples) & TN/FP/FN/TP reported in
notebook \\
\end{longtable}
}

ROC-AUC 0.87 indicates good ranking quality on a small dataset; accuracy
0.776 is typical for heart-disease prediction with these features and
reflects the problem's noise, not a modelling failure.

\section{8. Discussion and
Conclusion}\label{discussion-and-conclusion-1}

RFC gives a robust classifier with minimal preprocessing and
interpretable feature importances. On 227 training rows the
variance-reduction of bagging is important; class imbalance (about 45\%
positive here) is mild, but F1 and AUC are safer than accuracy alone.

\section{9. Answers to Viva
Questions}\label{answers-to-viva-questions-1}

\begin{itemize}
\tightlist
\item
  \textbf{RFR vs RFC?} Regression averages the trees' numeric outputs;
  classification takes a majority vote (and averages class probabilities
  for scoring).
\item
  \textbf{What is majority voting?} Each tree votes for one class and
  the most frequent class becomes the ensemble prediction (soft voting
  averages probabilities instead).
\item
  \textbf{How can class imbalance affect accuracy?} With a rare positive
  class, a model can predict the majority class everywhere and still
  show high accuracy; precision/recall/F1 and AUC expose the failure.
\end{itemize}

\begin{center}\rule{0.5\linewidth}{0.5pt}\end{center}

\section{10. Complete Code Listing}\label{complete-code-listing-1}

\emph{Full runnable program for this experiment, extracted from
\passthrough{\lstinline!notebooks/04\_ensemble\_learning.ipynb!}
(executed; results above are its actual output).}

\textbf{Listing 9.1 --- Core libraries}

\begin{lstlisting}[language=Python]
# ---- Core libraries ----
import numpy as np
import pandas as pd
import matplotlib.pyplot as plt
import seaborn as sns

from sklearn.datasets import fetch_california_housing
from sklearn.model_selection import train_test_split
from sklearn.ensemble import RandomForestRegressor, RandomForestClassifier, AdaBoostRegressor, AdaBoostClassifier
from sklearn.metrics import r2_score, root_mean_squared_error, mean_absolute_error, accuracy_score, f1_score, roc_auc_score

import xgboost as xgb
import catboost as cb

# Styling setup
plt.style.use('seaborn-v0_8-whitegrid' if 'seaborn-v0_8-whitegrid' in plt.style.available else 'default')
plt.rcParams['font.sans-serif'] = 'DejaVu Sans'
plt.rcParams['figure.figsize'] = (10, 6)
plt.rcParams['figure.dpi'] = 110
np.random.seed(42)
\end{lstlisting}

\textbf{Listing 9.2 --- Load the Heart Disease dataset (real-world
classification benchmark)}

\begin{lstlisting}[language=Python]
# ---- Load the Heart Disease dataset (real-world classification benchmark) ----
df_clf = pd.read_csv('../data/heart_disease.csv')
print(f"Heart Disease dataset: {df_clf.shape[0]} rows x {df_clf.shape[1]} columns")

# Separate target and features (handle either 'target' or 'Target' as the column name)
target_col = 'target' if 'target' in df_clf.columns else 'Target'
y_clf = df_clf[target_col].values
X_clf_raw = df_clf.drop(columns=[target_col])

# One-hot encode categorical variables for RFC / XGBoost / AdaBoost
X_clf_encoded = pd.get_dummies(X_clf_raw, drop_first=True)

# Stratified split keeps the disease/no-disease ratio identical in train and test
X_train_clf, X_test_clf, y_train_clf, y_test_clf = train_test_split(
    X_clf_encoded, y_clf, test_size=0.25, random_state=42, stratify=y_clf
)

print(f"Classification Training Set: {X_train_clf.shape}, Test Set: {X_test_clf.shape}")

# ---- Random Forest Classification (majority vote of trees) ----
rfc = RandomForestClassifier(n_estimators=120, max_depth=8, random_state=42, n_jobs=-1)
rfc.fit(X_train_clf, y_train_clf)

# Predictions and probability of the positive (disease) class
y_pred_rfc = rfc.predict(X_test_clf)
y_prob_rfc = rfc.predict_proba(X_test_clf)[:, 1]

# ---- Evaluate: accuracy, F1 and ROC-AUC ----
acc_rfc = accuracy_score(y_test_clf, y_pred_rfc)
f1_rfc = f1_score(y_test_clf, y_pred_rfc)
auc_rfc = roc_auc_score(y_test_clf, y_prob_rfc)

print(f"Random Forest Classifier: Accuracy = {acc_rfc:.4f}, F1 = {f1_rfc:.4f}, ROC-AUC = {auc_rfc:.4f}")
\end{lstlisting}

\chapter{Experiment 10: XGBoost}\label{experiment-10-xgboost}

\textbf{Category:} Ensemble learning / Gradient boosting

\section{1. Objective}\label{objective-2}

Train a gradient-boosted tree classifier and analyze the effect of
boosting hyperparameters.

\section{2. Dataset Used}\label{dataset-used-2}

Heart Disease (\passthrough{\lstinline!data/heart\_disease.csv!}), same
stratified 75/25 split as RFC for a fair comparison.

\section{3. Required Libraries}\label{required-libraries-2}

xgboost, scikit-learn, numpy, pandas

\section{4. Brief Theory}\label{brief-theory-2}

XGBoost builds trees sequentially, each fitting the 1st/2nd-order
gradients (g\_i, h\_i) of the loss: L ≈ Σ{[}g\_i f\_t(x\_i) + ½ h\_i
f\_t²(x\_i){]} + Ω(f\_t), with Ω(f) = γT + ½λΣw\_j². Controls:
\passthrough{\lstinline!learning\_rate!},
\passthrough{\lstinline!n\_estimators!},
\passthrough{\lstinline!max\_depth!},
\passthrough{\lstinline!subsample!},
\passthrough{\lstinline!colsample\_bytree!},
\passthrough{\lstinline!reg\_alpha!},
\passthrough{\lstinline!reg\_lambda!}.

\section{5. Procedure}\label{procedure-2}

\begin{enumerate}
\def\labelenumi{\arabic{enumi}.}
\tightlist
\item
  Prepare the one-hot encoded heart data and the same stratified split.
\item
  Fit
  \passthrough{\lstinline!XGBClassifier(n\_estimators=120, learning\_rate=0.08, max\_depth=4, subsample=0.8, colsample\_bytree=0.8, eval\_metric='logloss')!}.
\item
  Predict labels and probabilities.
\item
  Report accuracy, F1 and ROC-AUC.
\item
  Discuss the complexity/generalization trade-off.
\end{enumerate}

\section{6. Reference Implementation}\label{reference-implementation-2}

\begin{lstlisting}[language=Python]
xgb_clf = xgb.XGBClassifier(n_estimators=120, learning_rate=0.08, max_depth=4,
                            subsample=0.8, colsample_bytree=0.8,
                            random_state=42, eval_metric='logloss')
xgb_clf.fit(X_train_clf, y_train_clf)
y_prob_xgb_c = xgb_clf.predict_proba(X_test_clf)[:, 1]
\end{lstlisting}

\section{7. Results}\label{results-2}

{\def\LTcaptype{none} % do not increment counter
\begin{longtable}[]{@{}ll@{}}
\toprule\noalign{}
Metric & Value \\
\midrule\noalign{}
\endhead
\bottomrule\noalign{}
\endlastfoot
Accuracy & 0.7763 \\
F1 & 0.8132 \\
ROC-AUC & 0.8509 \\
\end{longtable}
}

XGBoost matches RFC on accuracy with a slightly higher F1 but a slightly
lower AUC on this small test set (76 samples). Differences of
\textasciitilde0.01-0.02 are within noise for this dataset size.

\section{8. Discussion and
Conclusion}\label{discussion-and-conclusion-2}

Gradient boosting is the strongest family for tabular data in general
(see regression results: XGBoost R² 0.8025 vs RFR 0.7415), but on this
small heart dataset with fixed hyperparameters it does not clearly beat
simpler ensembles. Lower learning rates with more trees usually improve
generalization at higher training cost.

\section{9. Answers to Viva
Questions}\label{answers-to-viva-questions-2}

\begin{itemize}
\tightlist
\item
  \textbf{Bagging vs boosting?} Bagging trains models in parallel on
  bootstrap samples and combines them to reduce variance; boosting
  trains models sequentially so each corrects the previous ensemble's
  errors, mainly reducing bias.
\item
  \textbf{What does learning rate do?} It shrinks each tree's
  contribution. Smaller values need more trees but usually generalize
  better; larger values fit faster and can overfit.
\item
  \textbf{Why can boosting overfit?} Later trees keep fitting remaining
  errors, including label noise, and the ensemble can become too
  complex; regularization (depth, λ, subsample, early stopping) controls
  this.
\end{itemize}

\begin{center}\rule{0.5\linewidth}{0.5pt}\end{center}

\section{10. Complete Code Listing}\label{complete-code-listing-2}

\emph{Full runnable program for this experiment, extracted from
\passthrough{\lstinline!notebooks/04\_ensemble\_learning.ipynb!}
(executed; results above are its actual output).}

\textbf{Listing 10.1 --- Core libraries}

\begin{lstlisting}[language=Python]
# ---- Core libraries ----
import numpy as np
import pandas as pd
import matplotlib.pyplot as plt
import seaborn as sns

from sklearn.datasets import fetch_california_housing
from sklearn.model_selection import train_test_split
from sklearn.ensemble import RandomForestRegressor, RandomForestClassifier, AdaBoostRegressor, AdaBoostClassifier
from sklearn.metrics import r2_score, root_mean_squared_error, mean_absolute_error, accuracy_score, f1_score, roc_auc_score

import xgboost as xgb
import catboost as cb

# Styling setup
plt.style.use('seaborn-v0_8-whitegrid' if 'seaborn-v0_8-whitegrid' in plt.style.available else 'default')
plt.rcParams['font.sans-serif'] = 'DejaVu Sans'
plt.rcParams['figure.figsize'] = (10, 6)
plt.rcParams['figure.dpi'] = 110
np.random.seed(42)
\end{lstlisting}

\textbf{Listing 10.2 --- Load the Heart Disease dataset (real-world
classification benchmark)}

\begin{lstlisting}[language=Python]
# ---- Load the Heart Disease dataset (real-world classification benchmark) ----
df_clf = pd.read_csv('../data/heart_disease.csv')
print(f"Heart Disease dataset: {df_clf.shape[0]} rows x {df_clf.shape[1]} columns")

# Separate target and features (handle either 'target' or 'Target' as the column name)
target_col = 'target' if 'target' in df_clf.columns else 'Target'
y_clf = df_clf[target_col].values
X_clf_raw = df_clf.drop(columns=[target_col])

# One-hot encode categorical variables for RFC / XGBoost / AdaBoost
X_clf_encoded = pd.get_dummies(X_clf_raw, drop_first=True)

# Stratified split keeps the disease/no-disease ratio identical in train and test
X_train_clf, X_test_clf, y_train_clf, y_test_clf = train_test_split(
    X_clf_encoded, y_clf, test_size=0.25, random_state=42, stratify=y_clf
)

print(f"Classification Training Set: {X_train_clf.shape}, Test Set: {X_test_clf.shape}")

# ---- Random Forest Classification (majority vote of trees) ----
rfc = RandomForestClassifier(n_estimators=120, max_depth=8, random_state=42, n_jobs=-1)
rfc.fit(X_train_clf, y_train_clf)

# Predictions and probability of the positive (disease) class
y_pred_rfc = rfc.predict(X_test_clf)
y_prob_rfc = rfc.predict_proba(X_test_clf)[:, 1]

# ---- Evaluate: accuracy, F1 and ROC-AUC ----
acc_rfc = accuracy_score(y_test_clf, y_pred_rfc)
f1_rfc = f1_score(y_test_clf, y_pred_rfc)
auc_rfc = roc_auc_score(y_test_clf, y_prob_rfc)

print(f"Random Forest Classifier: Accuracy = {acc_rfc:.4f}, F1 = {f1_rfc:.4f}, ROC-AUC = {auc_rfc:.4f}")
\end{lstlisting}

\textbf{Listing 10.3 --- XGBoost classification on the same stratified
split}

\begin{lstlisting}[language=Python]
# ---- XGBoost classification on the same stratified split ----
xgb_clf = xgb.XGBClassifier(
    n_estimators=120, learning_rate=0.08, max_depth=4,
    subsample=0.8, colsample_bytree=0.8,
    random_state=42, eval_metric='logloss'
)
xgb_clf.fit(X_train_clf, y_train_clf)

y_pred_xgb_c = xgb_clf.predict(X_test_clf)
y_prob_xgb_c = xgb_clf.predict_proba(X_test_clf)[:, 1]

acc_xgb_c = accuracy_score(y_test_clf, y_pred_xgb_c)
f1_xgb_c = f1_score(y_test_clf, y_pred_xgb_c)
auc_xgb_c = roc_auc_score(y_test_clf, y_prob_xgb_c)
print(f"XGBoost Classifier:   Accuracy = {acc_xgb_c:.4f}, F1 = {f1_xgb_c:.4f}, ROC-AUC = {auc_xgb_c:.4f}")
\end{lstlisting}

\chapter{Experiment 11: AdaBoost}\label{experiment-11-adaboost}

\textbf{Category:} Ensemble learning

\section{1. Objective}\label{objective-3}

Build an adaptive boosting classifier and study how weak learners
combine into a stronger ensemble.

\section{2. Dataset Used}\label{dataset-used-3}

Heart Disease (\passthrough{\lstinline!data/heart\_disease.csv!}), same
stratified 75/25 split.

\section{3. Required Libraries}\label{required-libraries-3}

scikit-learn, numpy, pandas

\section{4. Brief Theory}\label{brief-theory-3}

AdaBoost trains weak learners sequentially on re-weighted data:
misclassified samples receive larger weights. At round m the learner
weight is α\_m = ½ ln((1 − ε\_m)/ε\_m), and sample weights update as
w\_i ← w\_i exp(−α\_m y\_i h\_m(x\_i)). The final prediction is a
weighted vote. Exponential loss makes it sensitive to noisy points.

\section{5. Procedure}\label{procedure-3}

\begin{enumerate}
\def\labelenumi{\arabic{enumi}.}
\tightlist
\item
  Prepare the heart data and stratified split.
\item
  Fit
  \passthrough{\lstinline!AdaBoostClassifier(n\_estimators=100, learning\_rate=0.1)!}
  (stumps by default).
\item
  Predict labels and probabilities on the test set.
\item
  Report accuracy, F1 and ROC-AUC; compare with RFC and XGBoost.
\end{enumerate}

\section{6. Reference Implementation}\label{reference-implementation-3}

\begin{lstlisting}[language=Python]
ada_clf = AdaBoostClassifier(n_estimators=100, learning_rate=0.1, random_state=42)
ada_clf.fit(X_train_clf, y_train_clf)
y_prob_ada_c = ada_clf.predict_proba(X_test_clf)[:, 1]
\end{lstlisting}

\section{7. Results}\label{results-3}

{\def\LTcaptype{none} % do not increment counter
\begin{longtable}[]{@{}ll@{}}
\toprule\noalign{}
Metric & Value \\
\midrule\noalign{}
\endhead
\bottomrule\noalign{}
\endlastfoot
Accuracy & 0.7895 (best of the four classifiers) \\
F1 & 0.8182 (best) \\
ROC-AUC & 0.8714 (best) \\
\end{longtable}
}

AdaBoost with shallow stumps slightly outperforms the deeper models on
this small dataset, which is a common outcome when the signal is mostly
additive and the sample size is limited.

\section{8. Discussion and
Conclusion}\label{discussion-and-conclusion-3}

AdaBoost remains a strong, cheap baseline: stump ensembles are fast,
need little tuning, and here achieve top accuracy/F1/AUC. Its weakness
is sensitivity to label noise and outliers (weights grow exponentially
for hard points); gradient boosting with regularization is more robust
on noisier data.

\section{9. Answers to Viva
Questions}\label{answers-to-viva-questions-3}

\begin{itemize}
\tightlist
\item
  \textbf{Why use weak learners?} Individually they are simple and
  low-variance; boosting combines many of them so the ensemble can fit
  complex boundaries while keeping each step's variance low.
\item
  \textbf{How does AdaBoost focus on difficult examples?} After each
  round, weights of misclassified samples increase, so the next learner
  optimizes a distribution that emphasizes current errors.
\item
  \textbf{What is the effect of noisy labels?} Mislabeled points look
  ``hard'', receive ever-larger weights, and can dominate training,
  degrading accuracy --- AdaBoost is known to be noise-sensitive.
\end{itemize}

\begin{center}\rule{0.5\linewidth}{0.5pt}\end{center}

\section{10. Complete Code Listing}\label{complete-code-listing-3}

\emph{Full runnable program for this experiment, extracted from
\passthrough{\lstinline!notebooks/04\_ensemble\_learning.ipynb!}
(executed; results above are its actual output).}

\textbf{Listing 11.1 --- Core libraries}

\begin{lstlisting}[language=Python]
# ---- Core libraries ----
import numpy as np
import pandas as pd
import matplotlib.pyplot as plt
import seaborn as sns

from sklearn.datasets import fetch_california_housing
from sklearn.model_selection import train_test_split
from sklearn.ensemble import RandomForestRegressor, RandomForestClassifier, AdaBoostRegressor, AdaBoostClassifier
from sklearn.metrics import r2_score, root_mean_squared_error, mean_absolute_error, accuracy_score, f1_score, roc_auc_score

import xgboost as xgb
import catboost as cb

# Styling setup
plt.style.use('seaborn-v0_8-whitegrid' if 'seaborn-v0_8-whitegrid' in plt.style.available else 'default')
plt.rcParams['font.sans-serif'] = 'DejaVu Sans'
plt.rcParams['figure.figsize'] = (10, 6)
plt.rcParams['figure.dpi'] = 110
np.random.seed(42)
\end{lstlisting}

\textbf{Listing 11.2 --- Load the Heart Disease dataset (real-world
classification benchmark)}

\begin{lstlisting}[language=Python]
# ---- Load the Heart Disease dataset (real-world classification benchmark) ----
df_clf = pd.read_csv('../data/heart_disease.csv')
print(f"Heart Disease dataset: {df_clf.shape[0]} rows x {df_clf.shape[1]} columns")

# Separate target and features (handle either 'target' or 'Target' as the column name)
target_col = 'target' if 'target' in df_clf.columns else 'Target'
y_clf = df_clf[target_col].values
X_clf_raw = df_clf.drop(columns=[target_col])

# One-hot encode categorical variables for RFC / XGBoost / AdaBoost
X_clf_encoded = pd.get_dummies(X_clf_raw, drop_first=True)

# Stratified split keeps the disease/no-disease ratio identical in train and test
X_train_clf, X_test_clf, y_train_clf, y_test_clf = train_test_split(
    X_clf_encoded, y_clf, test_size=0.25, random_state=42, stratify=y_clf
)

print(f"Classification Training Set: {X_train_clf.shape}, Test Set: {X_test_clf.shape}")

# ---- Random Forest Classification (majority vote of trees) ----
rfc = RandomForestClassifier(n_estimators=120, max_depth=8, random_state=42, n_jobs=-1)
rfc.fit(X_train_clf, y_train_clf)

# Predictions and probability of the positive (disease) class
y_pred_rfc = rfc.predict(X_test_clf)
y_prob_rfc = rfc.predict_proba(X_test_clf)[:, 1]

# ---- Evaluate: accuracy, F1 and ROC-AUC ----
acc_rfc = accuracy_score(y_test_clf, y_pred_rfc)
f1_rfc = f1_score(y_test_clf, y_pred_rfc)
auc_rfc = roc_auc_score(y_test_clf, y_prob_rfc)

print(f"Random Forest Classifier: Accuracy = {acc_rfc:.4f}, F1 = {f1_rfc:.4f}, ROC-AUC = {auc_rfc:.4f}")
\end{lstlisting}

\textbf{Listing 11.3 --- AdaBoost classification}

\begin{lstlisting}[language=Python]
# ---- AdaBoost classification ----
ada_clf = AdaBoostClassifier(n_estimators=100, learning_rate=0.1, random_state=42)
ada_clf.fit(X_train_clf, y_train_clf)

y_pred_ada_c = ada_clf.predict(X_test_clf)
y_prob_ada_c = ada_clf.predict_proba(X_test_clf)[:, 1]

acc_ada_c = accuracy_score(y_test_clf, y_pred_ada_c)
f1_ada_c = f1_score(y_test_clf, y_pred_ada_c)
auc_ada_c = roc_auc_score(y_test_clf, y_prob_ada_c)
print(f"AdaBoost Classifier:  Accuracy = {acc_ada_c:.4f}, F1 = {f1_ada_c:.4f}, ROC-AUC = {auc_ada_c:.4f}")
\end{lstlisting}

\chapter{Experiment 12: CatBoost}\label{experiment-12-catboost}

\textbf{Category:} Ensemble learning

\section{1. Objective}\label{objective-4}

Train a gradient-boosted model with native handling of categorical
features and inspect correct categorical declaration.

\section{2. Dataset Used}\label{dataset-used-4}

Heart Disease (\passthrough{\lstinline!data/heart\_disease.csv!}):
continuous features (age, trestbps, chol, thalach, oldpeak) plus
categorical features (sex, cp, fbs, restecg, exang, slope, ca, thal).
The raw (unencoded) table is used for CatBoost.

\section{3. Required Libraries}\label{required-libraries-4}

catboost, pandas, scikit-learn

\section{4. Brief Theory}\label{brief-theory-4}

CatBoost uses ordered boosting (leaf values computed on random
permutations) to avoid target leakage/prediction shift, and builds
oblivious (symmetric) trees. Categorical features are encoded internally
with ordered target statistics, so
\passthrough{\lstinline!cat\_features!} must be declared correctly and
the data split must not leak the target.

\section{5. Procedure}\label{procedure-4}

\begin{enumerate}
\def\labelenumi{\arabic{enumi}.}
\tightlist
\item
  Identify categorical columns in the raw heart table and cast them to
  string.
\item
  Use the same stratified 75/25 split on the raw features.
\item
  Fit
  \passthrough{\lstinline!CatBoostClassifier(iterations=150, depth=5, cat\_features=categorical\_cols)!}.
\item
  Predict labels/probabilities; report accuracy, F1, ROC-AUC.
\item
  Explain why target-statistic encoding can beat one-hot encoding.
\end{enumerate}

\section{6. Reference Implementation}\label{reference-implementation-4}

\begin{lstlisting}[language=Python]
categorical_cols = X_clf_raw.select_dtypes(include=['object', 'category']).columns.tolist()
cb_clf = cb.CatBoostClassifier(iterations=150, learning_rate=0.08, depth=5,
                               cat_features=categorical_cols, random_seed=42, verbose=False)
cb_clf.fit(X_tr_cb, y_tr_cb)
\end{lstlisting}

\section{7. Results}\label{results-4}

{\def\LTcaptype{none} % do not increment counter
\begin{longtable}[]{@{}ll@{}}
\toprule\noalign{}
Metric & Value \\
\midrule\noalign{}
\endhead
\bottomrule\noalign{}
\endlastfoot
Accuracy & 0.7632 \\
F1 & 0.7955 \\
ROC-AUC & 0.8551 \\
\end{longtable}
}

CatBoost performs at a similar level to the other ensembles without any
manual one-hot encoding --- its native categorical handling did the
encoding internally. On this tiny dataset, the advanced target
statistics cannot show their full advantage (they matter most with
high-cardinality categories, e.g.~thousands of city/user IDs).

\section{8. Discussion and
Conclusion}\label{discussion-and-conclusion-4}

CatBoost is the ensemble of choice when tables mix numeric and
categorical columns: it removes manual encoding, reduces target leakage,
and remains strong out-of-the-box. Here its metrics sit within noise of
RFC/XGBoost/AdaBoost because the categorical cardinalities are small
(2-4 values).

\section{9. Answers to Viva
Questions}\label{answers-to-viva-questions-4}

\begin{itemize}
\tightlist
\item
  \textbf{How does CatBoost treat categorical features?} It computes
  ordered target statistics (mean target per category on random
  permutations with prior smoothing) rather than one-hot vectors.
\item
  \textbf{Why should categories not be encoded using target leakage?}
  Plain target encoding uses each row's own target when encoding that
  row, which leaks label information and inflates training performance;
  ordered statistics avoid this.
\item
  \textbf{CatBoost vs ordinary one-hot encoding?} One-hot explodes
  dimensionality for high-cardinality features and loses category
  statistics; CatBoost encodes compactly and exploits target-category
  relationships with leakage control.
\end{itemize}

\begin{center}\rule{0.5\linewidth}{0.5pt}\end{center}

\section{10. Complete Code Listing}\label{complete-code-listing-4}

\emph{Full runnable program for this experiment, extracted from
\passthrough{\lstinline!notebooks/04\_ensemble\_learning.ipynb!}
(executed; results above are its actual output).}

\textbf{Listing 12.1 --- Core libraries}

\begin{lstlisting}[language=Python]
# ---- Core libraries ----
import numpy as np
import pandas as pd
import matplotlib.pyplot as plt
import seaborn as sns

from sklearn.datasets import fetch_california_housing
from sklearn.model_selection import train_test_split
from sklearn.ensemble import RandomForestRegressor, RandomForestClassifier, AdaBoostRegressor, AdaBoostClassifier
from sklearn.metrics import r2_score, root_mean_squared_error, mean_absolute_error, accuracy_score, f1_score, roc_auc_score

import xgboost as xgb
import catboost as cb

# Styling setup
plt.style.use('seaborn-v0_8-whitegrid' if 'seaborn-v0_8-whitegrid' in plt.style.available else 'default')
plt.rcParams['font.sans-serif'] = 'DejaVu Sans'
plt.rcParams['figure.figsize'] = (10, 6)
plt.rcParams['figure.dpi'] = 110
np.random.seed(42)
\end{lstlisting}

\textbf{Listing 12.2 --- Load the Heart Disease dataset (real-world
classification benchmark)}

\begin{lstlisting}[language=Python]
# ---- Load the Heart Disease dataset (real-world classification benchmark) ----
df_clf = pd.read_csv('../data/heart_disease.csv')
print(f"Heart Disease dataset: {df_clf.shape[0]} rows x {df_clf.shape[1]} columns")

# Separate target and features (handle either 'target' or 'Target' as the column name)
target_col = 'target' if 'target' in df_clf.columns else 'Target'
y_clf = df_clf[target_col].values
X_clf_raw = df_clf.drop(columns=[target_col])

# One-hot encode categorical variables for RFC / XGBoost / AdaBoost
X_clf_encoded = pd.get_dummies(X_clf_raw, drop_first=True)

# Stratified split keeps the disease/no-disease ratio identical in train and test
X_train_clf, X_test_clf, y_train_clf, y_test_clf = train_test_split(
    X_clf_encoded, y_clf, test_size=0.25, random_state=42, stratify=y_clf
)

print(f"Classification Training Set: {X_train_clf.shape}, Test Set: {X_test_clf.shape}")

# ---- Random Forest Classification (majority vote of trees) ----
rfc = RandomForestClassifier(n_estimators=120, max_depth=8, random_state=42, n_jobs=-1)
rfc.fit(X_train_clf, y_train_clf)

# Predictions and probability of the positive (disease) class
y_pred_rfc = rfc.predict(X_test_clf)
y_prob_rfc = rfc.predict_proba(X_test_clf)[:, 1]

# ---- Evaluate: accuracy, F1 and ROC-AUC ----
acc_rfc = accuracy_score(y_test_clf, y_pred_rfc)
f1_rfc = f1_score(y_test_clf, y_pred_rfc)
auc_rfc = roc_auc_score(y_test_clf, y_prob_rfc)

print(f"Random Forest Classifier: Accuracy = {acc_rfc:.4f}, F1 = {f1_rfc:.4f}, ROC-AUC = {auc_rfc:.4f}")
\end{lstlisting}

\textbf{Listing 12.3 --- CatBoost classification with native categorical
features}

\begin{lstlisting}[language=Python]
# ---- CatBoost classification with native categorical features ----
# Identify categorical columns in the raw (unencoded) heart data
categorical_cols = X_clf_raw.select_dtypes(include=['object', 'category']).columns.tolist()
X_raw_cb = X_clf_raw.copy()
for col in categorical_cols:
    X_raw_cb[col] = X_raw_cb[col].astype(str)

# Same stratified split on the raw feature table
X_tr_cb, X_te_cb, y_tr_cb, y_te_cb = train_test_split(
    X_raw_cb, y_clf, test_size=0.25, random_state=42, stratify=y_clf
)

cb_clf = cb.CatBoostClassifier(
    iterations=150, learning_rate=0.08, depth=5,
    cat_features=categorical_cols, random_seed=42, verbose=False
)
cb_clf.fit(X_tr_cb, y_tr_cb)

y_pred_cb_c = cb_clf.predict(X_te_cb)
y_prob_cb_c = cb_clf.predict_proba(X_te_cb)[:, 1]

acc_cb_c = accuracy_score(y_te_cb, y_pred_cb_c)
f1_cb_c = f1_score(y_te_cb, y_pred_cb_c)
auc_cb_c = roc_auc_score(y_te_cb, y_prob_cb_c)
print(f"CatBoost Classifier:  Accuracy = {acc_cb_c:.4f}, F1 = {f1_cb_c:.4f}, ROC-AUC = {auc_cb_c:.4f}")
\end{lstlisting}
